%BAB IX: KESIMPULAN DAN SARAN
\section{Kesimpulan dan Saran}

\subsection{Kesimpulan}

Kegiatan magang di PT Sumber Cipta Multiniaga (GDP Labs) selama periode 9 Februari 2026 -- 22 Juni 2026 memberikan kontribusi signifikan terhadap pengembangan keterampilan teknis dan non-teknis penulis. Dari sisi hardskill, pengalaman ini meliputi:

\begin{enumerate}
    \item \textbf{Pengembangan fitur dan modul}: Merancang dan mengimplementasikan custom tools untuk Digital Employee (tenant info tool, weekly report tools, GitHub issue automation) dengan prinsip modularitas dan clean code.
    \item \textbf{Implementasi prototipe}: Membangun proof-of-concept untuk skenario AI seperti issue scoring di E2B sandbox, pipeline RAG (semantic routing, query transformation), serta analisis masukan AI pada pull request.
    \item \textbf{Pengembangan backend}: Mengembangkan layanan backend dengan Python dan FastAPI yang terintegrasi dengan GL-IAM SDK untuk autentikasi berbasis JWT, manajemen API key, dan agent delegation.
    \item \textbf{Pengujian unit dan modul}: Menyusun dan menjalankan pengujian unit menggunakan pytest dengan pendekatan mocking dan analisis cakupan kode (coverage).
    \item \textbf{Pengujian integrasi dan sistem}: Mengikuti serta berkontribusi pada pengujian integrasi melalui CI/CD (GitHub Actions) dan pengujian end-to-end dengan Playwright untuk memastikan keandalan sistem secara menyeluruh.
\end{enumerate}

Dari sisi softskill, kegiatan ini melatih kemampuan bekerjasama dalam tim Agile/Scrum, komunikasi teknis melalui code review, serta manajemen tugas menggunakan GitHub Projects.

\subsection{Saran}

Berdasarkan pengalaman mengikuti program MBKM di GDP Labs, penulis menyampaikan saran berikut:

\subsubsection{Saran untuk Perusahaan}
\begin{enumerate}
    \item \textbf{Dokumentasi runbook untuk CI/CD failures}: Tim infra disarankan untuk membuat runbook resmi yang mendeskripsikan daftar kegagalan pipeline umum dan langkah penanganannya, sehingga mengurangi waktu investigasi perulangan.
    \item \textbf{Peningkatan onboarding proyek baru}: Template Cookiecutter sudah sangat membantu, namun penambahan panduan video atau sesi live walkthrough akan mempercepat adaptasi anggota tim baru, termasuk magang.
\end{enumerate}

\subsubsection{Saran untuk Akademik}
\begin{enumerate}
    \item \textbf{Integrasi praktik Agile ke kurikulum}: Disarankan agar mata kuliah Proyek Rekayasa Perangkat Lunak atau sejenisnya memasukkan elemen praktik Agile (sprint planning, stand-up, retrospektif) sebagai bagian dari penilaian proyek, sehingga mahasiswa lebih siap menghadapi lingkungan industri.
    \item \textbf{Pembelajaran infrastructure as code}: Menambahkan topik deployment (Docker, Kubernetes, Helm) ke kurikulum opsional akan sangat bermanfaat mengingat banyak perusahaan teknologi menggunakan pendekatan ini untuk deployment produksi.
\end{enumerate}

\clearpage

